\documentclass[a4paper]{article}

\usepackage[a4paper, top=8mm, bottom=8mm, left=8mm, right=8mm]{geometry}

\usepackage{polyglossia}
\setdefaultlanguage[babelshorthands=true]{russian}
\setotherlanguage{english}

\usepackage{fontspec}
\setmainfont{FreeSerif}
\newfontfamily{\russianfonttt}[Scale=0.7]{DejaVuSansMono}

\usepackage[tiny, compact]{titlesec}

\usepackage{titling}
\setlength{\droptitle}{-1cm}
\pretitle{\begin{center}\begin{bfseries}\Large}
\posttitle{\par\end{bfseries}\end{center}}
\preauthor{\begin{center}\normalsize}
\postauthor{\par\end{center}\vspace{-1.8cm}}

\usepackage{hyperref}
\usepackage{bookmark}
\usepackage{csquotes}

\title{UCFS: Доработка репозитория}

\author{Исаев Даниил Сергеевич}

\date{}

\begin{document}

\maketitle

\begin{flushright}
    Группа:~\emph{25.Б72-мм}

    Кафедра:~\emph{Кафедра системного программирования СПбГУ}

    Научный руководитель:~\emph{Григорьев Семён Вячеславович}
    
    Номер семестра практики:~\emph{2}
\end{flushright} 

\section{Введение}

Работа выполняется в рамках проекта UCFS\footnote{Ссылка на UCFS: \url{https://github.com/FormalLanguageConstrainedPathQuerying/UCFS}}. UCFS~--- это универсальный инструмент, предназначенный для решения задач, которые лежат на пересечении контекстно-свободных языков и ориентированных графов с метками на ребрах. В основе его работы лежит алгоритм GLL~(Generalized LL parsing)~--- алгоритм синтаксического анализа, обобщающий LL-анализ на случай неоднозначных грамматик. Этот инструмент применим для широкого круга задач, включая синтаксический анализ и контекстно-свободную навигацию по графам (CFPQ/CFL-R).

На данный момент проект обладает широкой функциональностью, однако существующие примеры использования сложны, так как решают конкретную узкую задачу и не позволяют новичкам разобраться в использовании инструмента. Данная практика направлена на улучшение документации и создание пошаговой инструкции использования инструмента, которая завершается запуском простых примеров, что позволит быстрее освоить инструмент и начать применять его в учебных и исследовательских задачах.

\section{Постановка задачи}

\textbf{Целью работы является} создание пошаговой инструкции по аналогии с популярными библиотеками, а также простых примеров использования инструмента UCFS, включая реализацию кода и интеграцию инструкции по запуску в основную документацию.

\textbf{Для достижения цели требуется решить следующие задачи:}
\begin{enumerate}
    \item Изучить материалы касающиеся контекстно-свободных грамматик и языков.
    \item Выполнить обзор узкой задачи, решаемой с использованием инструмента UCFS, и подходов к написанию документации в популярных проектах.
    \item Сделать вывод о том, какой информации не хватает для быстрого старта использования UCFS новичками и спроектировать минимальную, но информативную инструкцию.
    \item Спроектировать набор минимальных учебных примеров и формат их интеграции в документацию.
    \item Реализовать разработанные примеры на языке Kotlin с использованием встроенного DSL и подготовить визуализации графов.
    \item Выполнить тестирование реализованных примеров на корректность работы.
    \item Выполнить реструктуризацию документации проекта:~оставить README файлы техническими, всю побочную информацию вынести в документацию.
\end{enumerate}

\section{Обзор}

В рамках проекта UCFS уже существовал набор примеров использования. Однако анализ текущей документации показал, что существующие примеры оперируют сложными грамматиками и большими графами, взятыми из реальных прикладных задач. Это затрудняет понимание базовых принципов работы инструмента пользователями, которые только начинают знакомство с CFPQ и GLL-анализом.

В других проектах, таких как NetworkX\footnote{Ссылка на документацию NetworkX: \url{https://networkx.org/documentation/stable/}}, стандартной практикой является наличие разделов \enquote{Tutorial}, в котором содержится инструкция по работе с инструментом, и \enquote{Getting Started} с минимальными воспроизводимыми примерами, которые демонстрируют базовый API. В данном проекте вышеперечисленные разделы отсутствуют.

\textbf{Вывод:} Таким образом, в текущем виде документация UCFS не подходит для быстрого старта, и возникает необходимость в создании упрощенных учебных сценариев с пошаговой инструкцией по запуску.

\section{Описание предлагаемого решения}

Предлагаемое решение заключается в создании документации, подробно поясняющей каждый шаг новичка для начала работы с UCFS, а именно базовую информацию о CFL~(Context-Free Languages), инструкцию по заданию графов и грамматик, инструкцию по чтению выходных данных, а также содержащую набор минимальных и понятных примеров, которые позволяют последовательно познакомиться с возможностями UCFS. В отличие от существовавших ранее сложных примеров, разработанные сценарии намеренно упрощены:~графы невелики по размеру, а их структура интуитивно понятна. Ключевым примером выбрана грамматика $a^nb^n$~--- она не имеет принципиального значения для демонстрации возможностей инструмента, но была взята из учебного пособия\footnote{Ссылка на пособие: \url{https://github.com/FormalLanguageConstrainedPathQuerying/FormalLanguageConstrainedReachability-LectureNotes} (дата обращения:~14.03.2026).}, что обеспечивает методическую обоснованность выбора. Для описания грамматики использовался встроенный DSL проекта UCFS, а графы представлены в формате~\texttt{.dot}.

Принятые решения обоснованы следующим образом:~простота примеров снижает порог входа и подготавливает новых пользователей к работе со сложными графами, имеющими практическое применение. Интеграция инструкции по запуску непосредственно в основной файл документации (вместе с визуализациями графов в формате~SVG)~устраняет необходимость поиска отдельной информации.

Реализация решения осуществлялась с использованием языка Kotlin~(код примеров), DSL для задания грамматики, формата~\texttt{.dot} для описания графов, консольной утилиты dot~(Graphviz)~для конвертации графов в SVG и Markdown для оформления документации.

\section{Тестирование}

Проверка разработанных примеров осуществлялась посредством тестирования. Каждый граф был вручную исследован на корректность структуры, после чего с помощью инструмента UCFS для него был получен SPPF~(Shared Packed Parse Forest~--- общий упакованный лес разбора). Анализ полученного SPPF показал, что результаты совпадают с ожидаемыми результатами поиска по графу, что подтверждает корректность работы примеров.

Запуск примеров выполняется единой командой через консольный сценарий. Проверка переносимости показала, что проблем с запуском не возникает благодаря стандартизированному формату графов~(\texttt{.dot}) и отсутствию внешних зависимостей, специфичных для окружения разработки.

Оценка успешности работы проводилась по следующим критериям:~все примеры успешно компилируются и запускаются, результаты их работы (построенные леса разбора) совпадают с теоретически ожидаемыми для выбранной грамматики и графов.

\section{Заключение}

В ходе выполнения работы получены следующие результаты:

\begin{itemize}
    \item Разработаны простые примеры использования инструмента UCFS.
    \item Созданы графы в формате~\texttt{.dot} для наглядной демонстрации работы инструмента.
    \item Разработанные примеры прошли проверку через анализ SPPF, подтверждающую корректность их работы.
    \item Разработаны новые разделы документации для использования UCFS, готовые к интеграции в основной файл.
    \item Создан пуллреквест, который на данный момент находится на рассмотрении.
\end{itemize}

Материалы, иллюстрирующие выполнение работы, размещены в пуллреквесте.

Репозиторий (форк):~\url{https://github.com/danisaev/UCFS}.

Ссылка на пуллреквест:~\url{https://github.com/vkutuev/UCFS/pull/3}.

После успешного прохождения ревью разработанные примеры будут включены в основную ветку проекта UCFS в качестве учебных материалов для новых пользователей.

\end{document}
